\documentclass[11pt]{article}
\usepackage{ctex}
\usepackage{geometry}
\usepackage{hyperref}
\usepackage{listings}
\usepackage{xcolor}
\geometry{a4paper, margin=1in}

\lstset{
	basicstyle=\ttfamily\small,
	breaklines=true,
	columns=fullflexible,
	frame=single,
	rulecolor=\color{black},
	xleftmargin=1em,
	xrightmargin=1em
}

\title{项目修改记录（1\_31）}
\author{}
\date{2026-02-01}

\begin{document}
\maketitle

\section*{概览}
1.保留之前构建数据集的脚本build\_dataset，并把生成csv改成生成train.py需要的Pkl。\par
2.统一格式把build\_dataset里的参数放到options.py里。\par
3.在train.py里修改接受pkl的逻辑。\par
4.由于train.py里面的cnn处理的timing report这些信息我们没有，所以去掉了CNN/path相关逻辑。\par
5.参考之前的train\_hgat将train.py里的gnn替换成HGAT。\par
6.在options.py里修改train相关参数，比如特征维度对齐。\par

\section{build\_dataset.py}
\subsection*{修改 1：统一从 options 读取参数}
\textbf{修改前：}
\begin{lstlisting}
parser = argparse.ArgumentParser()
parser.add_argument("--src_lib", required=True, ...)
...
args = parser.parse_args()
\end{lstlisting}

\textbf{修改后：}
\begin{lstlisting}
from options import get_options
...
args = get_options()
\end{lstlisting}

\textbf{说明：} 统一参数入口，便于后续集中管理与复现实验配置。

\subsection*{修改 2：生成 dataset.pkl（训练直接读取）}
\textbf{修改前：}
\begin{lstlisting}
# 仅输出 CSV 和 json
\end{lstlisting}

\textbf{修改后：}
\begin{lstlisting}
dataset_obj = {
    "meta": meta,
    "scaler_stats": scaler_stats,
    "y_scaler": y_scaler,
    "tgt_train_df": df_tgt_train,
    "tgt_val_df": df_tgt_val,
    "tgt_test_df": df_tgt_test,
    "src_df": df_src,
}
with open(os.path.join(args.out_dir, args.dataset_pkl_name), "wb") as f:
    pickle.dump(dataset_obj, f)
\end{lstlisting}

\textbf{说明：} 训练脚本以 pkl 作为统一数据接口，避免依赖 CSV 中间文件。

\subsection*{修改 3：数据切分参数化}
\textbf{修改前：}
\begin{lstlisting}
train_cells, val_cells, test_cells = split_by_cell_type(df_tgt,
    ratios=(0.14, 0.14, 0.72), seed=42)
\end{lstlisting}

\textbf{修改后：}
\begin{lstlisting}
train_cells, val_cells, test_cells = split_by_cell_type(
    df_tgt,
    ratios=tuple(args.tgt_split_ratios),
    seed=args.split_seed,
)
\end{lstlisting}

\textbf{说明：} 切分比例与随机种子可通过 options 配置，便于实验对比。

\subsection*{修改 4：Slew 阈值归一化（跨工艺对齐）}
\textbf{修改前：}
\begin{lstlisting}
# 直接使用 lib 中的 slew 数值
row = _build_enhanced_row(..., slew=float(s), ...)
\end{lstlisting}

\textbf{修改后：}
\begin{lstlisting}
thresholds = parse_slew_thresholds(lib_text)
span = (upper - lower) / 100.0
slew_norm = s * derate / span
row = _build_enhanced_row(..., slew=slew_norm, ...)
\end{lstlisting}

\textbf{说明：} 不同工艺库采用不同阈值定义（如 10--90\% 与 30--70\%）。
为避免跨工艺的尺度偏差，在构建阶段按阈值区间严格归一化。

\section{options.py}
\subsection*{修改 4：新增 build\_dataset 配置项}
\textbf{修改前：}
\begin{lstlisting}
--src_lib --tgt_lib --src_spi --tgt_sp --out_dir --target_label_ratio
\end{lstlisting}

\textbf{修改后：}
\begin{lstlisting}
--tgt_split_ratios 0.14 0.14 0.72
--split_seed 42
--dataset_pkl_name dataset.pkl
\end{lstlisting}

\textbf{说明：} 将切分比例、随机种子与 pkl 文件名纳入统一参数管理。

\subsection*{修改 5：训练默认数据集路径对齐}
\textbf{修改前：}
\begin{lstlisting}
--data_save_path ../datasets/asap7-designs
\end{lstlisting}

\textbf{修改后：}
\begin{lstlisting}
--data_save_path ../output
\end{lstlisting}


\section{train.py}
\subsection*{修改 6：支持可配置的 pkl 名称}
\textbf{修改前：}
\begin{lstlisting}
load_dataset_pkl(data_dir)  # 固定 dataset.pkl
\end{lstlisting}

\textbf{修改后：}
\begin{lstlisting}
load_dataset_pkl(data_dir, options.dataset_pkl_name)
# 若不存在则回退到 dataset.pkl
\end{lstlisting}

\textbf{说明：} 训练读取与 build\_dataset 输出的 pkl 文件名保持一致。

\subsection*{修改 7：删除 CNN/path/timing report 逻辑}
\textbf{修改前：}
\begin{lstlisting}
   if val_loss < best_val:
...
th.save({...}, ...)

# endpoints = ...
sampled_ends, sampled_paths = {}, {}
...
if options.no_cnn or len(paths) == 0:
path_map = None
else:
path_map = path_mask.to_dense() * feat_map
...
val_res, val_F1_score, val_r2 = validate(...)
...
\end{lstlisting}

\textbf{修改后：}
\begin{lstlisting}
if val_loss < best_val:
...
th.save({...}, ...)
\end{lstlisting}

\textbf{说明：} 当前数据源仅包含 lib 与 sp，无 timing report 与 path 信息，保留这些逻辑会导致运行错误。

\subsection*{修改8：模型初始化（PathConv+CNN+MLP → HGAT+MLP）}
\textbf{修改前：}
\begin{lstlisting}
if options.no_gnn:
gnn = None
else:
gnn = PathConv(
in_feats_dim=options.out_dim,
out_feats_dim=options.out_dim,
cell_feat_dim=options.cell_feat_dim,
net_feat_dim=options.net_feat_dim,
flag_attn=options.attn,
num_heads=options.num_heads
)
mlp_dim += gnn.out_feats_dim

if options.no_cnn:
cnn = None
fcn = None
else:
cnn = UNet(options.pooling) if options.unet else LayoutNet(options.pooling)
fcn = nn.Linear(options.map_size * options.map_size, options.cnn_outdim)
mlp_dim += options.cnn_outdim

mlp = MLP2(in_dim=mlp_dim, out_dim=options.nlabels, nlayers=options.n_fcn, dropout=options.mlp_dropout)
model = PathModel(gnn, fcn, mlp)
\end{lstlisting}

\textbf{修改后：}
\begin{lstlisting}
in_map = {"NET": 4, "PMOS": 2, "NMOS": 2}
enc = HGATDesignEncoder(in_dim_map=in_map, hid=hgat_hid, out=design_dim, num_heads=hgat_heads).to(device)
model = CellDelayRegressor(in_dim=options.in_dim, design_dim=design_dim, hid=hgat_hid, dropout=dropout).to(device)
\end{lstlisting}


\subsection*{修改9：训练输入（路径/布局特征 → 预计算 HGAT cell embedding）}
\textbf{修改前：}
\begin{lstlisting}
for path_dataset, graph, path2level, path2endpoint, topo_levels, cnn_inputs, path_masks in loader:
cnn_inputs = cnn_inputs.reshape((1, cnn_inputs.shape[0], cnn_inputs.shape[1], cnn_inputs.shape[2]))
feat_map = cnn(cnn_inputs.to(device)).reshape((1, -1)) if cnn is not None else None
graph = graph.to(device)
...
\end{lstlisting}

\textbf{修改后：}
\begin{lstlisting}
print("----------------Loading HGAT embeddings----------------")
z_dict = precompute_z_from_tgt_spice(data_dir, meta, enc, device, design_dim)

def z_provider(ct):
z = z_dict.get(ct)
if z is None:
return th.zeros(1, design_dim, device=device)
return z

for xb, yb, cts in train_dl:
xb = xb.to(device)
yb = yb.to(device)
zb = th.cat([z_provider(ct) for ct in cts], dim=0)
pred = model(xb, zb)
\end{lstlisting}

\subsection*{修改10：训练输入特征维度校验}
\textbf{修改前：}
\begin{lstlisting}
model = CellDelayRegressor(in_dim=len(NUMERIC_COLS), ...)
\end{lstlisting}

\textbf{修改后：}
\begin{lstlisting}
if options.in_dim != len(NUMERIC_COLS):
    raise ValueError(...)
model = CellDelayRegressor(in_dim=options.in_dim, ...)
\end{lstlisting}



\end{document}
